\medskip

\begin{minipage}{0.58\linewidth}
On considère le prisme droit ABFEDCGH, de base ABFE, trapèze rectangle en A. 

On associe à ce prisme le repère orthonormé $\left(\text{A}~;~\vect{\imath}, \vect{\jmath},\vect{k}\right)$ tel que : 

\[\vect{\imath} = \dfrac14\vect{\text{AB}}, \quad \vect{\jmath} = \dfrac14\vect{\text{AD}}, \quad \vect{k} = \dfrac18\vect{\text{AE}}.\]

De plus on a $\vect{\text{BF}} = \dfrac12\vect{\text{AE}}$.

On note I le milieu du segment [EF].

On note J le milieu du segment [AE].
\end{minipage}\hfill
\begin{minipage}{0.35\linewidth}
\psset{unit=1cm}
\begin{pspicture}(6,8)

\psline[linewidth=1.25pt]{->}(0.2,1.4)(5.7,0)
\psline[linewidth=1.25pt]{->}(0.2,1.4)(0.2,7.3)
\psline[linewidth=1.25pt,linestyle=dashed](0.2,1.4)(4.2,3.2)
\psline[linewidth=1.25pt]{->}(4.2,3.2)(5.2,3.7)
\psline(2.6,0.8)(4.2,1.6)(4.2,4.2)(2.6,3.4)(2.6,0.8)
\psline(4.2,4.2)(1.8,7.3)(0.2,6.8)(2.6,3.4)
\psline(0.2,4)(1.35,5.2)(4.2,4.2)
\psline(1.35,5.2)(1.8,7.3)
\psline(2.6,0.8)(4.2,4.2)
\psline[linestyle=dashed](4.2,1.6)(1.8,2.1)(1.8,7.3)
\psline[linestyle=dashed](1.8,7.3)(0.2,4)(4.2,4.2)
\psline[linewidth=1.25pt]{->}(0.2,1.4)(0.8,1.24)
\psline[linewidth=1.25pt]{->}(0.2,1.4)(0.7,1.6)
\psline[linewidth=1.25pt]{->}(0.2,1.4)(0.2,2.2)
\uput[dl](0.2,1.4){A} \uput[d](2.6,0.8){B} \uput[r](4.2,1.6){C} \uput[ul](1.8,2.1){D}
\uput[l](0.2,6.8){E} \uput[dl](2.6,3.4){F} \uput[r](4.2,4.2){G} \uput[u](1.8,7.3){H}
\uput[d](1.35,5.2){I} \uput[l](0.2,4){J} \uput[d](0.5,1.3){$\vect{\imath}$} \uput[ur](0.35,1.5){$\vect{\jmath}$}
\uput[l](0.2,1.8){$\vect{k}$}
\end{pspicture}
\end{minipage}
\bigskip

\begin{enumerate}
\item Donner les coordonnées des points I et J.
\item Soit $\vect{n}$ le vecteur de coordonnées $\begin{pmatrix}-1\\1\\1\end{pmatrix}$.
	\begin{enumerate}
		\item Montrer que le vecteur $\vect{n}$ est normal au plan (IGJ).
		\item Déterminer une équation cartésienne du plan (IGJ).
	\end{enumerate}	
\item Déterminer une représentation paramétrique de la droite $d$, perpendiculaire au
plan (IGJ) et passant par H.
\item On note L le projeté orthogonal du point H sur le plan (IGJ).

Montrer que les coordonnées de L sont $\left(\dfrac83~;~ \dfrac43~;~\dfrac{16}{3}\right)$.
\item Calculer la distance du point H au plan (IGJ).
\item Montrer que le triangle IGJ est rectangle en I.
\item En déduire le volume du tétraèdre IGJH.
\end{enumerate}

\begin{center}
On rappelle que le volume $V$ d'un tétraèdre est donné par la formule :

$V = \dfrac13 \times (\text{aire de la base}) \times  \text{hauteur}.$
\end{center}

\bigskip
